\clearpage
\VerdierPhysicalPageAnchor{0249}%
\begingroup\small\setlength{\parskip}{0pt}\setlength{\emergencystretch}{1.5em}
\setlength{\abovedisplayskip}{4pt plus 2pt minus 2pt}%
\setlength{\belowdisplayskip}{4pt plus 2pt minus 2pt}%

\VerdierSemanticAnchor{ch4d-4-6-9-first-differentials}
\noindent et les premières différentielles sont (4.1.9.1) :
\[
\begin{aligned}
d_1^{p,q}\colon(\mathrm{C.E.})_1^{p,q}
  &\longrightarrow(\mathrm{C.E.})_1^{p+1,q},\\
d_1^{p,q}&=F^q(-d_2^{\bullet,p}).
\end{aligned}
\]

\VerdierSemanticAnchor{ch4d-4-6-9-column-decomposition}
Les complexes $I^{\bullet,p}$ sont des complexes injectifs de
Cartan--Eilenberg (4.6.1) et, par suite (4.6.3), on a un isomorphisme
canonique dans $D(\mathcal A)$ :
\[
I^{\bullet,p}\xrightarrow{\sim}
\prod_{r\in\mathbb Z}H^r(I^{\bullet,p})[-r].
\]

\VerdierSemanticAnchor{ch4d-4-6-9-h1-components}
Utilisons alors le complexe $H_1^r(I^{\bullet\bullet})$ introduit en
(4.6.4.3). Pour tout entier $r$ on a :
\[
H_1^r(I^{\bullet\bullet})^p=H^r(I^{\bullet,p}),
\]
d'où un isomorphisme canonique dans $D(\mathcal A)$ :

\VerdierSemanticAnchor{ch4d-4-6-9-h1-decomposition}
\[
I^{\bullet,p}\xrightarrow{\sim}
\prod_{r\in\mathbb Z}H_1^r(I^{\bullet\bullet})^p[-r],
\]
et, modulo cet isomorphisme, le morphisme de complexes :
\[
d_2^{\bullet,p}\colon
I^{\bullet,p}\longrightarrow I^{\bullet,p+1}
\]
est isomorphe au morphisme :

\VerdierSemanticAnchor{ch4d-4-6-9-d2-decomposition}
\[
\prod_{r\in\mathbb Z}
d_{H_1^r(I^{\bullet\bullet})}^{p}[-r],
\]
d'où l'expression définitive des termes initiaux :

\VerdierSemanticAnchor{ch4d-4-6-9-1}
\[
(\mathrm{C.E.})_1^{p,q}
=F^q\left(
\prod_{r\in\mathbb Z}H_1^r(I^{\bullet\bullet})^p[-r]
\right),
\tag{4.6.9.1}
\]
et des premières différentielles :

\VerdierSemanticAnchor{ch4d-4-6-9-2}
\[
d_1^{p,q}
=-F^q\left(
\prod_{r\in\mathbb Z}d_{H_1^r(I^{\bullet\bullet})}^{p}[-r]
\right).
\tag{4.6.9.2}
\]

\VerdierSemanticAnchor{ch4d-4-6-9-3-prelude}
Soit alors $F^{q,\bullet}(H_1^r(I^{\bullet\bullet}))$ le complexe d'objets
de $\mathcal B$ obtenu en appliquant le foncteur $F^q$ aux composants et
aux différentielles du complexe $H_1^r(I^{\bullet\bullet})$. Les formules
(4.6.9.1) et (4.6.9.2) définissent un morphisme :

\VerdierSemanticAnchor{ch4d-4-6-9-3}
\[
(\mathrm{C.E.})_2^{p,q}\longrightarrow
\prod_{r\in\mathbb Z}
H_{\mathcal B}^{p}\left(
F^{q-r,\bullet}(H_1^r(I^{\bullet\bullet}))
\right).
\tag{4.6.9.3}
\]
\endgroup
